\documentclass[10pt,a4paper]{amsart}
\usepackage[margin=19mm]{geometry}
\usepackage{amsmath,amssymb,mathtools}
\usepackage[scr=rsfs,cal=euler]{mathalfa}
\usepackage{xfrac}
\usepackage[unicode,hidelinks,bookmarks=false]{hyperref}
\newcommand{\ie}{i.e.\ }
\newcommand{\cf}{cf.\ }
\newcommand{\resp}{respectively\ }
\newcommand{\Subsection}{\subsection}
\providecommand{\nobreakdash}{\nobreak\hbox{-}}
\setlength{\emergencystretch}{2em}
\multlinegap0pt
\usepackage{amssymb}
\usepackage{mathtools}
\newtheorem{corollary}{Corollary}
\newtheorem{proposition}{Proposition}
\newcommand{\norm}[1]{\Vert#1\Vert}
\title{Positive isometric Fourier multipliers on non-commutative $L^p$-spaces}
\author{Christoph Kriegler, Christian Le Merdy, Safoura Zadeh}
\date{Source: arXiv:2603.07754v1 (CC BY 4.0)}
\begin{document}
\maketitle

\noindent\emph{Source and licence.} Positive isometric Fourier multipliers on non-commutative $L^p$-spaces. Version arXiv:2603.07754v1 (CC BY 4.0), distributed by arXiv under the Creative Commons Attribution 4.0 International licence, http://creativecommons.org/licenses/by/4.0/, as recorded in the arXiv licence metadata for this exact version and shown on the abstract page https://arxiv.org/abs/2603.07754v1. Full licence notice: \texttt{licenses/arxiv-2603.07754-CC-BY-4.0.txt}.\par\smallskip

\noindent\textit{Editorial scope.} Counted unit: the article's corollary 5SelfDual (source0.tex lines 1123--1129). The article's complete original proof of it is delivered with it (source0.tex lines 1132--1143, one \texttt{proof} environment of 12 lines). No mathematics is written, repaired or re-derived: the statement and the proof are the article's own lines, in the article's own notation, and no retained line is substituted or re-flowed.  Retained uncounted as the source-local context the proof needs: lines 635--638; lines 744--752; lines 1068--1072.  Nothing was omitted from any retained span, so the package carries no omission record.  No retained line contains a citation, so the wrapper carries no bibliography and no bibliography entry.  No retained line contains a figure environment, an image-inclusion command or drawing code, so no image is delivered and the package declares no asset dependency.  Editorial formatting only, in addition: an AMS \texttt{amsart} preamble that restates the article's own declarations and macros used by the retained lines, namely the environments \texttt{corollary}, \texttt{proposition} on separate counters, together with the standard packages those lines need.  The retained environments print in this document as Corollary 1, Proposition 1 in order of appearance, whereas the article prints the same items with the numbering of its own full document.  The complete native source of the article as distributed in the arXiv e-print for this exact version is bundled unchanged as \texttt{sources/195961/source0.tex}.
\begin{corollary}\label{3Always}
Let $\phi\in L^\infty(G)$. Then $\phi$ is a bounded Fourier multiplier on $L^2({\mathcal L}G)$ and moreover,
$\norm{M_{\phi,2}}= \norm{\phi}_\infty.$
\end{corollary}

\begin{proposition}\label{3Duality} 
Let $1<p<\infty$, let $p'$ denote its conjugate number and let $\phi\in L^\infty(G)$. If $\phi$
is a bounded Fourier multiplier on $L^p({\mathcal L}G)$, 
then $\check{\phi}$ is a bounded Fourier multiplier on $L^{p'}({\mathcal L}G)$,
and 
$
M_{\phi,p}^*= M_{\check{\phi},p'}.
$
\end{proposition}

\begin{proposition}\label{4Antipode}
Let $\phi\in L^\infty(G)$ and let $1\leq p<\infty$. If $\phi$ is a bounded Fourier multiplier on
$L^p({\mathcal L}G)$, then $\check{\phi}$ is also a bounded Fourier multiplier, with
$\norm{M_{\check{\phi},p}}= \norm{M_{\phi,p}}$.
\end{proposition}

\begin{corollary}\label{5SelfDual}
Let $\phi\in L^\infty(G)$, let $1< p<\infty$ and let $p'$ be its conjugate number. 
Then $\phi$ is a bounded Fourier multiplier on
$L^p({\mathcal L}G)$ if and only if $\phi$ is a bounded Fourier multiplier on
$L^{p'}({\mathcal L}G)$. In this case, we
have $\norm{M_{\phi,p'}}=\norm{M_{\phi,p}}$ and $\norm{\phi}_\infty\leq \norm{M_{\phi,p}}$.
\end{corollary}

\begin{proof}
Assume that $\phi\in L^\infty(G)$
is a bounded Fourier multiplier on
$L^p({\mathcal L}G)$.
Combining Proposition \ref{4Antipode} and Proposition \ref{3Duality}, we obtain
that $\phi$ is a bounded Fourier multiplier on
$L^{p'}({\mathcal L}G)$, with $\norm{M_{\phi,p'}}=\norm{M_{\phi,p}}$. 
Then by interpolation,
$\norm{M_{\phi,2}}\leq \norm{M_{\phi,p}}$ and hence 
$\norm{\phi}_\infty\leq \norm{M_{\phi,p}}$
by Corollary \ref{3Always}.
\end{proof}

\end{document}
